\documentclass[12pt,a4paper]{article}
\usepackage[utf8]{inputenc}
\usepackage{amsmath, amssymb}
\usepackage{enumitem}
\usepackage[margin=2cm]{geometry}

\begin{document}

\begin{center}
    \textbf{\Large Latihan Soal Asesmen Sumatif Akhir Tahun (ASAT)} \\
    \textbf{Bahasa Indonesia Kelas X --- SMK Plus Pratama Adi}
\end{center}
\vspace{0.5cm}

\begin{enumerate}[leftmargin=*]

% --- TEKS LAPORAN HASIL OBSERVASI (1-15) ---

\item Tujuan utama dari penulisan teks laporan hasil observasi (LHO) adalah \dots
    \begin{enumerate}[label=\alph*.]
        \item Menghibur pembaca dengan cerita fiksi hasil imajinasi penulis.
        \item Mengajak pembaca untuk mengikuti opini pribadi penulis.
        \item Menyajikan informasi faktual mengenai suatu objek sesuai hasil pengamatan.
        \item Membujuk masyarakat untuk membeli suatu produk hasil penelitian.
        \item Mengkritik kebijakan pemerintah secara tidak langsung.
    \end{enumerate}

\item Teks laporan hasil observasi ditulis untuk memenuhi kebutuhan berikut, kecuali \dots
    \begin{enumerate}[label=\alph*.]
        \item Memberikan informasi berbasis data yang valid.
        \item Mengatasi suatu masalah atau persoalan secara objektif.
        \item Mendokumentasikan hasil pengamatan ilmiah.
        \item Mengungkapkan perasaan emosional penulis secara subjektif.
        \item Mengetahui perkembangan dan kemajuan suatu objek.
    \end{enumerate}

\item Susunan urutan struktur teks laporan hasil observasi yang benar adalah \dots
    \begin{enumerate}[label=\alph*.]
        \item Deskripsi bagian, pernyataan umum, deskripsi manfaat.
        \item Pernyataan umum, deskripsi bagian, deskripsi manfaat.
        \item Deskripsi manfaat, pernyataan umum, deskripsi bagian.
        \item Pernyataan umum, deskripsi manfaat, deskripsi bagian.
        \item Deskripsi bagian, deskripsi manfaat, kesimpulan umum.
    \end{enumerate}

\item Urutan struktur teks LHO yang ideal diawali dengan pembukaan yang memuat \dots
    \begin{enumerate}[label=\alph*.]
        \item Penggolongan atau definisi dasar mengenai objek yang diamati.
        \item Rincian manfaat fungsional dari objek dalam kehidupan sehari-hari.
        \item Penjelasan detail mengenai anatomi atau bagian fisik objek.
        \item Analisis mendalam mengenai dampak negatif dari keberadaan objek.
        \item Kesimpulan akhir yang merangkum hasil evaluasi menyeluruh.
    \end{enumerate}

\item Di bawah ini yang merupakan ciri kebahasaan dominan dari teks laporan hasil observasi adalah \dots
    \begin{enumerate}[label=\alph*.]
        \item Banyak menggunakan kata kerja ekspresif dan kata ganti orang pertama.
        \item Menggunakan kalimat perintah, seruan, dan bahasa kiasan yang padat.
        \item Menggunakan kalimat definisi, verba relasional, dan kata ilmiah/teknis.
        \item Menggunakan konjungsi temporal kronologis dan kalimat tanya retoris.
        \item Didominasi oleh majas sarkasme dan penggunaan ungkapan sindiran.
    \end{enumerate}

\item Perhatikan kata-kata berikut: \textit{ekosistem, mamalia, fotosintesis, habitat}. Kata-kata tersebut termasuk ciri kebahasaan teks LHO berupa \dots
    \begin{enumerate}[label=\alph*.]
        \item Kata sandang
        \item Istilah teknis/ilmiah
        \item Kata kerja material
        \item Frasa konotatif
        \item Kata ganti penunjuk
    \end{enumerate}

\item Cermati kutipan teks berikut:\\
\textit{"Kantong semar memiliki daun berbentuk piala yang berfungsi menangkap serangga. Di dalam piala tersebut terdapat cairan asam encer yang berfungsi mencerna tubuh serangga."}\\
Berdasarkan isinya, kutipan tersebut termasuk bagian \dots
    \begin{enumerate}[label=\alph*.]
        \item Pernyataan umum
        \item Klasifikasi objek
        \item Deskripsi bagian
        \item Deskripsi manfaat
        \item Penegasan ulang
    \end{enumerate}

\item Cermati kutipan teks berikut:\\
\textit{"Kunang-kunang adalah sejenis serangga yang dapat mengeluarkan cahaya yang jelas terlihat pada malam hari."}\\
Kutipan di atas merupakan bagian struktur LHO yang dinamakan \dots
    \begin{enumerate}[label=\alph*.]
        \item Deskripsi bagian
        \item Pernyataan umum
        \item Deskripsi manfaat
        \item Kesimpulan ringkas
        \item Saran penulis
    \end{enumerate}

\item Kalimat yang sesuai dengan karakteristik teks laporan hasil observasi harus bersifat \dots
    \begin{enumerate}[label=\alph*.]
        \item Imajinatif dan persuasif
        \item Subjektif dan argumentatif
        \item Objektif dan faktual
        \item Puitis dan ekspresif
        \item Fiktif dan solutif
    \end{enumerate}

\item Perhatikan teks berikut:\\
\textit{"Hutan bakau memiliki sistem perakaran yang rapat. Akar-akar tersebut mampu menahan hantaman ombak laut sehingga mencegah terjadinya abrasi di sepanjang garis pantai."}\\
Informasi utama yang terdapat dalam kutipan observasi tersebut adalah \dots
    \begin{enumerate}[label=\alph*.]
        \item Hutan bakau banyak tumbuh di wilayah daratan tinggi.
        \item Akar hutan bakau sangat rentan terhadap erosi air sungai.
        \item Fungsi akar hutan bakau dalam mencegah abrasi pantai.
        \item Keindahan hutan bakau sebagai tempat wisata air.
        \item Cara membudidayakan tanaman bakau di pesisir pantai.
    \end{enumerate}

\item Di antara kalimat-kalimat berikut, yang tidak objektif untuk digunakan dalam teks LHO adalah \dots
    \begin{enumerate}[label=\alph*.]
        \item Komodo memiliki panjang tubuh rata-rata mencapai 2 hingga 3 meter.
        \item Bunga melati putih itu terlihat sangat cantik dan menawan hati siapa saja.
        \item Tanaman lidah buaya tumbuh subur di daerah yang beriklim tropis.
        \item Paus pembunuh sebenarnya termasuk dalam keluarga lumba-lumba besar.
        \item Kelelawar merupakan satu-satunya mamalia yang memiliki kemampuan terbang.
    \end{enumerate}

\item Perhatikan hubungan isi antarparagraf berikut:\\
Paragraf pertama menjelaskan klasifikasi jenis-jenis sampah, sedangkan paragraf kedua menjelaskan cara pengolahan sampah organik menjadi kompos. Hubungan isi kedua paragraf tersebut adalah \dots
    \begin{enumerate}[label=\alph*.]
        \item Hubungan sebab dan akibat yang berlawanan.
        \item Hubungan pengenalan umum yang dilanjutkan rincian spesifik objek.
        \item Hubungan pertentangan ide pokok antartokoh.
        \item Hubungan ruang dan latar waktu masa lampau.
        \item Hubungan analogi fiktif antardua benda.
    \end{enumerate}

\item Pernyataan berikut ini yang tidak mencerminkan prinsip objektivitas teks laporan hasil observasi adalah \dots
    \begin{enumerate}[label=\alph*.]
        \item Berdasarkan data laboratorium, kadar air dalam madu asli berkisar 17-20\%.
        \item Saya merasa bahwa kucing persia adalah hewan yang paling malas di dunia.
        \item Candi Borobudur dibangun menggunakan batuan andesit yang disusun rapi.
        \item Harimau sumatra merupakan salah satu satwa yang terancam punah.
        \item Lumba-lumba berkomunikasi menggunakan sistem sonar frekuensi tinggi.
    \end{enumerate}

\item Kalimat berikut ini yang memerlukan perbaikan agar menjadi kalimat efektif dan baku untuk teks LHO adalah \dots
    \begin{enumerate}[label=\alph*.]
        \item Banyak hewan-hewan liar mati karena kelaparan di musim kemarau.
        \item Beruang kutub memiliki lapisan lemak yang tebal di bawah kulitnya.
        \item Kebun raya berfungsi sebagai pusat konservasi tumbuhan ex-situ.
        \item Terumbu karang membutuhkan sinar matahari untuk berfotosintesis.
        \item Landak mempertahankan diri dengan menegakkan duri pada tubuhnya.
    \end{enumerate}

\item Kalimat yang mengalami pemborosan kata dan perlu diperbaiki adalah \dots
    \begin{enumerate}[label=\alph*.]
        \item Tanaman itu disiram setiap pagi hari agar tidak layu.
        \item Menurut para ahli ekologi menyatakan bahwa hutan adalah paru-paru dunia.
        \item Mangrove merupakan pelindung alami kawasan pesisir pantai.
        \item Satwa langka dilindungi oleh undang-undang yang berlaku.
        \item Semut bekerja sama untuk mengumpulkan cadangan makanan mereka.
    \end{enumerate}


% --- TEKS ANEKDOT (16-30) ---

\item Tujuan utama dari penulisan sebuah teks anekdot adalah \dots
    \begin{enumerate}[label=\alph*.]
        \item Menyampaikan petunjuk teknis pengoperasian alat.
        \item Menyampaikan kritik atau sindiran sosial melalui cerita lucu.
        \item Memaparkan keindahan alam suatu daerah wisata.
        \item Menjelaskan riwayat hidup seorang tokoh fiktif secara detail.
        \item Membujuk pembaca agar menyetujui program politik tertentu.
    \end{enumerate}

\item Bagian dari struktur teks anekdot yang memuat masalah utama, kejanggalan, atau puncak kelucuan cerita disebut \dots
    \begin{enumerate}[label=\alph*.]
        \item Abstraksi
        \item Orientasi
        \item Krisis
        \item Reaksi
        \item Koda
    \end{enumerate}

\item Di bawah ini yang merupakan ciri kebahasaan teks anekdot dari segi penggunaan waktu adalah \dots
    \begin{enumerate}[label=\alph*.]
        \item Menggunakan kata kerja masa depan yang akan datang.
        \item Menggunakan keterangan waktu lampau yang sudah terjadi.
        \item Menggunakan kalimat pasif berimbuhan ter- secara dominan.
        \item Hanya menggunakan istilah-istilah medis dan biologi ilmiah.
        \item Tidak menggunakan konjungsi temporal sama sekali.
    \end{enumerate}

\item Cermati kutipan teks anekdot berikut:\\
\textit{"Seorang pejabat korup mendatangi seorang peramal untuk menanyakan nasibnya setelah masa jabatannya berakhir tahun depan. Peramal itu pun tersenyum tipis."}\\
Bagian kutipan tersebut di dalam struktur anekdot termasuk bagian \dots
    \begin{enumerate}[label=\alph*.]
        \item Abstraksi
        \item Orientasi
        \item Krisis
        \item Reaksi
        \item Koda
    \end{enumerate}

\item Elemen humor di dalam sebuah teks anekdot ditandai oleh hal berikut, kecuali \dots
    \begin{enumerate}[label=\alph*.]
        \item Adanya situasi janggal yang dialami tokoh.
        \item Tanggapan atau jawaban tokoh yang tidak terduga/absurd.
        \item Sindiran menggelitik yang memicu senyuman pembaca.
        \item Penjelasan rumus ilmiah yang kaku dan matematis.
        \item Permainan kata atau kesalahpahaman komunikasi antartokoh.
    \end{enumerate}

\item Kelucuan atau penyebab munculnya humor dalam suatu cerita anekdot biasanya berasal dari \dots
    \begin{enumerate}[label=\alph*.]
        \item Keadaan sedih tokoh yang mengalami bencana alam dahsyat.
        \item Ketidakselarasan logis atau \textit{plot twist} di akhir dialog tokoh.
        \item Penggunaan bahasa asing yang sangat baku dan formal.
        \item Deskripsi detail mengenai latar tempat yang mewah dan megah.
        \item Banyaknya tokoh pembantu yang tidak melakukan aksi apa pun.
    \end{enumerate}

\item Cermati kutipan dialog anekdot berikut:\\
\textit{A: "Mengapa kamu suka membuang sampah sembarangan di sungai?"}\\
\textit{B: "Biar sungai kita terlihat memiliki banyak isi, tidak kosong seperti janji kampanye bapak."}\\
Kritik sosial yang ingin disampaikan lewat kutipan di atas berkaitan dengan \dots
    \begin{enumerate}[label=\alph*.]
        \item Buruknya fasilitas pendidikan di daerah pedalaman.
        \item Kurangnya kesadaran warga dalam menjaga kesehatan tubuh.
        \item Sindiran terhadap janji politik pejabat yang tidak ditepati.
        \item Masalah pencemaran udara akibat gas emisi kendaraan.
        \item Kenaikan harga bahan pokok menjelang hari raya.
    \end{enumerate}

\item Kritik dalam teks anekdot harus disampaikan secara santun, artinya \dots
    \begin{enumerate}[label=\alph*.]
        \item Menggunakan kata-kata kasar agar langsung dipahami sasaran.
        \item Dikemas dalam bentuk cerita lucu atau sindiran tersirat yang halus.
        \item Ditulis dalam bentuk laporan ilmiah resmi yang bertele-tele.
        \item Menyerang fisik atau privasi tokoh yang dikritik secara vulgar.
        \item Menggunakan bahasa isyarat yang tidak dipahami oleh umum.
    \end{enumerate}

\item Perhatikan kalimat berikut:\\
\textit{"Pak Hakim memvonis pencuri sandal 5 tahun penjara, sedangkan koruptor milyaran rupiah hanya divonis 1 tahun penjara."}\\
Sasaran sindiran dari kalimat anekdot tersebut ditujukan kepada \dots
    \begin{enumerate}[label=\alph*.]
        \item Warga masyarakat yang gemar memakai sandal jepit.
        \item Ketidakadilan dan ketimpangan penegakan hukum di suatu negara.
        \item Para pembuat undang-undang hak cipta barang mewah.
        \item Penjual sandal yang mengalami kerugian materiil besar.
        \item Petugas kebersihan yang lalai menjaga keamanan kantor.
    \end{enumerate}

\item Efektivitas kritik sosial dalam sebuah teks anekdot dianggap berhasil apabila \dots
    \begin{enumerate}[label=\alph*.]
        \item Pembaca merasa tersinggung dan marah setelah membaca cerita.
        \item Pembaca tertawa sekaligus tersadar akan pesan moral/sindiran di dalamnya.
        \item Cerita tersebut ditulis sepanjang sepuluh halaman penuh tanpa jeda.
        \item Menggunakan tokoh hewan fiktif yang tidak bisa berbicara.
        \item Tidak ada satu pun kalimat dialog yang digunakan dalam teks.
    \end{enumerate}

\item Evaluasi ketepatan isi terhadap sebuah pesan teks anekdot harus didasarkan pada \dots
    \begin{enumerate}[label=\alph*.]
        \item Seberapa seram latar tempat yang digambarkan penulis.
        \item Relevansi isu sosial yang dikritik dengan realitas kehidupan nyata.
        \item Keindahan rima dan bait sajak akhir yang digunakan.
        \item Kebenaran rumus-rumus fisika yang dimasukkan ke dalam teks.
        \item Banyaknya majas hiperbola yang membesar-besarkan fakta fiksi.
    \end{enumerate}

\item Teks anekdot berbeda dengan teks humor biasa. Perbedaan esensial tersebut terletak pada \dots
    \begin{enumerate}[label=\alph*.]
        \item Teks humor selalu panjang, sedangkan anekdot pasti pendek.
        \item Teks humor memuat kritik sosial, sedangkan anekdot hanya menghibur.
        \item Teks anekdot memuat kritik sosial, sedangkan humor biasa hanya untuk menghibur.
        \item Teks anekdot menggunakan bahasa ilmiah, sedangkan humor menggunakan bahasa puitis.
        \item Teks humor selalu menggunakan tokoh nyata, sedangkan anekdot menggunakan alien.
    \end{enumerate}

\item Kalimat tanya retoris yang sering ditemukan dalam teks anekdot memiliki karakteristik \dots
    \begin{enumerate}[label=\alph*.]
        \item Membutuhkan jawaban ilmiah yang panjang dari para ahli.
        \item Pertanyaan yang jawabannya sudah jelas dan tidak perlu dijawab.
        \item Digunakan untuk menanyakan arah jalan yang tersesat.
        \item Memerlukan data statistik kuantitatif sebelum dijawab.
        \item Hanya boleh diucapkan oleh tokoh utama yang berprofesi sebagai hakim.
    \end{enumerate}

\item Perhatikan teks anekdot acak berikut:\\
1) Petugas kebersihan itu pun kebingungan.\\
2) Seorang dosen mengkritik kebijakan kampus di depan mahasiswanya.\\
3) Mahasiswa tertawa mendengar analogi cerdas sang dosen.\\
4) Di akhir kuliah, sang dosen terpeleset kulit pisang.\\
Bagian krisis pada rangkaian cerita di atas ditunjukkan ketika \dots
    \begin{enumerate}[label=\alph*.]
        \item Dosen mengkritik kebijakan kampus dengan analogi unik.
        \item Mahasiswa tertawa bersama-sama di ruang kelas.
        \item Dosen terpeleset kulit pisang akibat kelalaian bersama.
        \item Petugas kebersihan kebingungan menyapu lantai kampus.
        \item Mahasiswa pulang ke rumah masing-masing setelah selesai kuliah.
    \end{enumerate}

\item Jika sebuah anekdot mengkritik perilaku masyarakat yang gemar membuang puntung rokok sembarangan, maka sasaran sindiran utamanya adalah \dots
    \begin{enumerate}[label=\alph*.]
        \item Pabrik pembuat korek api gas.
        \item Kurangnya kesadaran masyarakat terhadap kebersihan lingkungan.
        \item Petugas medis yang mengobati penyakit paru-paru.
        \item Para menteri yang mengatur impor daun tembakau.
        \item Penjual asbak hias di pasar malam tradisional.
    \end{enumerate}


% --- PUISI (31-45) ---

\item Jenis puisi lama yang tiap baitnya terdiri dari empat baris, bersajak a-b-a-b, baris pertama dan kedua merupakan sampiran, sedangkan baris ketiga dan keempat merupakan isi dinamakan \dots
    \begin{enumerate}[label=\alph*.]
        \item Syair
        \item Gurindam
        \item Pantun
        \item Karmina
        \item Talibun
    \end{enumerate}

\item Puisi lama yang terdiri atas dua baris dalam satu bait, bersajak a-a, dan biasanya berisi nasihat atau filosofi hidup dinamakan \dots
    \begin{enumerate}[label=\alph*.]
        \item Syair
        \item Gurindam
        \item Pantun
        \item Seloka
        \item Soneta
    \end{enumerate}

\item Perhatikan ciri-ciri puisi berikut:\\
1) Setiap bait terdiri atas 4 baris.\\
2) Seluruh baris merupakan bagian isi.\\
3) Memiliki rima akhir bersajak a-a-a-a.\\
Berdasarkan ciri-ciri di atas, jenis puisi lama yang dimaksud adalah \dots
    \begin{enumerate}[label=\alph*.]
        \item Pantun
        \item Gurindam
        \item Syair
        \item Karmina
        \item Distikhon
    \end{enumerate}

\item Musikalisasi puisi adalah sebuah konsep karya seni yang dilakukan dengan cara \dots
    \begin{enumerate}[label=\alph*.]
        \item Mengubah seluruh lirik lagu pop menjadi bait-bait puisi lama.
        \item Membaca puisi yang dikolaborasikan atau diiringi dengan alunan musik.
        \item Menulis puisi di atas lembaran kertas partitur not balok musik.
        \item Menyanyikan lagu tanpa menghayati makna kata-kata di dalamnya.
        \item Menghapus teks puisi dan menggantinya dengan instrumen gitar murni.
    \end{enumerate}

\item Cermati kutipan puisi berikut:\\
\textit{Wahai pemuda kelana cita}\\
\textit{Janganlah bimbang menatap dunia}\\
\textit{Singsingkan lengan bulatkan karsa}\\
\textit{Demi masa depan indonesia}\\
Berdasarkan isinya, kutipan di atas digolongkan sebagai puisi jenis \dots
    \begin{enumerate}[label=\alph*.]
        \item Balada
        \item Elegi
        \item Syair/Nasihat Pemuda
        \item Romansa
        \item Satire
    \end{enumerate}

\item Aspek dominan yang harus dinilai saat seseorang membaca puisi di atas panggung (deklamasi) adalah \dots
    \begin{enumerate}[label=\alph*.]
        \item Kecepatan berlari dan kekuatan fisik deklamator.
        \item Vokal (artikulasi, intonasi), ekspresi wajah, dan gestur tubuh.
        \item Keindahan pakaian adat dan tata rias yang digunakan pembaca.
        \item Merk mikrofon dan sistem pengeras suara yang dipasang panitia.
        \item Jumlah penonton yang bersorak riuh di dalam ruangan gedung.
    \end{enumerate}

\item Cermati kutipan puisi berikut:\\
\textit{Ibu, di keheningan malam yang sunyi}\\
\textit{Doamu mengalir bagai air jernih tanpa henti}\\
\textit{Membasuh luka hati yang perih ini}\\
Suasana dominan yang tergambar dalam kutipan puisi tersebut adalah \dots
    \begin{enumerate}[label=\alph*.]
        \item Riang gembira dan penuh tawa
        \item Syahdu, penuh haru, dan khidmat
        \item Tegang, marah, dan menakutkan
        \item Panik dan terburu-buru
        \item Lucu dan menggelitik sanubari
    \end{enumerate}

\item Cermati kutipan puisi karya W.S. Rendra berikut:\\
\textit{Kita adalah sejuta tangan di jalanan}\\
\textit{Yang menuntut hak atas sejuta piring nasi}\\
\textit{Yang dirampas oleh keserakahan tirani}\\
Kritik sosial yang terdapat dalam kutipan puisi tersebut menyuarakan tentang \dots
    \begin{enumerate}[label=\alph*.]
        \item Keindahan alam desa yang mulai luntur akibat industri.
        \item Protes terhadap ketidakadilan ekonomi dan keserakahan penguasa.
        \item Kegembiraan para buruh yang menerima bonus tahunan.
        \item Ajakan untuk menanam padi di sawah milik bersama.
        \item Keluh kesah petani akibat musim hujan yang datang terlambat.
    \end{enumerate}

\item Di dalam puisi, kata \textit{"senja"} sering kali digunakan sebagai lambang atau makna simbolik untuk menyatakan \dots
    \begin{enumerate}[label=\alph*.]
        \item Masa kanak-kanak yang penuh permainan.
        \item Harapan baru yang muncul di pagi hari.
        \item Masa tua atau fase akhir perjalanan kehidupan manusia.
        \item Kekayaan materiil yang berlimpah ruah.
        \item Semangat perjuangan pemuda yang berkobar.
    \end{enumerate}

\item Cermati kutipan puisi berikut:\\
\textit{Jangan biarkan mimpimu padam oleh waktu}\\
\textit{Teruslah melangkah lurus ke depan}\\
\textit{Meski kerikil tajam melukai tapak kakimu}\\
Amanat yang ingin disampaikan oleh penyair lewat kutipan puisi tersebut adalah \dots
    \begin{enumerate}[label=\alph*.]
        \item Kita harus segera membeli sepatu baru agar kaki tidak terluka.
        \item Jangan pernah bermimpi di siang hari karena membuang waktu.
        \item Kita harus tetap tegar dan pantang menyerah dalam mengejar cita-cita.
        \item Sebaiknya hindari jalan yang dipenuhi oleh batu kerikil tajam.
        \item Mintalah bantuan orang lain saat menghadapi masalah hidup yang sulit.
    \end{enumerate}

\item Penilaian terhadap penggunaan diksi dalam puisi dilakukan dengan cara melihat \dots
    \begin{enumerate}[label=\alph*.]
        \item Jumlah total kata serapan dari bahasa asing yang digunakan penulis.
        \item Ketepatan pemilihan kata yang mampu menimbulkan efek estetika dan kedalaman makna.
        \item Kesesuaian ejaan kata dengan kamus besar bahasa Indonesia versi terbaru.
        \item Panjang atau pendeknya rangkaian huruf dalam satu baris kalimat.
        \item Jumlah tinta yang dihabiskan untuk mencetak baris-baris puisi.
    \end{enumerate}

\item Ciri utama dari puisi modern (puisi bebas) yang membedakannya dengan puisi lama adalah \dots
    \begin{enumerate}[label=\alph*.]
        \item Tidak terikat oleh aturan ketat mengenai rima, bait, dan jumlah suku kata.
        \item Wajib menggunakan bahasa daerah kuno yang sulit dipahami awam.
        \item Hanya boleh ditulis oleh para sastrawan senior yang bergelar doktor.
        \item Selalu mengisahkan tentang kehidupan para dewa kuno di khayangan.
        \item Harus diiringi oleh instrumen musik gamelan jawa tradisional.
    \end{enumerate}

\item Jika sebuah puisi menggunakan kata \textit{"cemara yang gugur sebatang demi sebatang"}, suasana yang ingin dibangun oleh penyair adalah \dots
    \begin{enumerate}[label=\alph*.]
        \item Semangat juang membara
        \item Kesunyian, kesedihan, atau kedukaan
        \item Kepanikan akibat bencana perang
        \item Kebahagiaan pesta pernikahan
        \item Ketakutan di dalam hutan belantara
    \end{enumerate}

\item Kata \textit{"samudra"} dalam puisi sering kali dipilih sebagai lambang yang merepresentasikan \dots
    \begin{enumerate}[label=\alph*.]
        \item Kesempitan ruang lingkup berpikir manusia.
        \item Keluasan hati, ilmu pengetahuan, atau perjuangan hidup yang luas.
        \item Sifat manusia yang cepat marah dan emosional.
        \item Kekeringan yang melanda ladang pertanian warga.
        \item Rasa haus yang mendalam setelah berolahraga siang hari.
    \end{enumerate}

\item Cermati kutipan puisi singkat berikut:\\
\textit{Hiduplah seperti lilin kecil yang rela terbakar demi menerangi kegelapan sekitarnya.}\\
Amanat yang tersirat dari kutipan di atas adalah \dots
    \begin{enumerate}[label=\alph*.]
        \item Ajakan untuk selalu menghemat penggunaan listrik rumah tangga.
        \item Kita harus memiliki sifat rela berkorban demi memberikan kebaikan bagi sesama.
        \item Larangan menyalakan api di dalam ruangan tertutup tanpa pengawasan.
        \item Saran untuk memproduksi lilin hias sebagai komoditas ekspor.
        \item Kita harus menjadi orang yang lemah dan mudah hancur dihantam cobaan.
    \end{enumerate}


% --- TEKS EKSPOSISI & AFIKSASI (46-60) ---

\item Tujuan utama dari penulisan teks eksposisi adalah \dots
    \begin{enumerate}[label=\alph*.]
        \item Menceritakan urutan peristiwa fiksi secara kronologis objektif.
        \item Memaparkan informasi atau pengetahuan tertentu disertai argumen logis agar wawasan pembaca bertambah.
        \item Menggambarkan objek secara mendetail hingga pembaca seolah melihatnya sendiri.
        \item Melukiskan perasaan sedih penyair melalui untaian kata indah berpola.
        \item Memandu pembaca melakukan langkah eksperimen zat kimia berbahaya.
    \end{enumerate}

\item Ciri utama yang paling menonjol dari sebuah teks eksposisi adalah informasinya bersifat \dots
    \begin{enumerate}[label=\alph*.]
        \item Subjektif, penuh mitos, dan imajinatif
        \item Ilmiah, logis, faktual, dan tidak memihak
        \item Khayalan, puitis, dan menggunakan rima akhir
        \item Rahasia, tertutup, dan hanya dibaca kelompok tertentu
        \item Naratif fiktif dengan penekanan pada tokoh antagonis
    \end{enumerate}

\item Bagian struktur teks eksposisi yang memuat gagasan utama, pandangan penulis, atau opini awal mengenai suatu topik permasalahan disebut \dots
    \begin{enumerate}[label=\alph*.]
        \item Tesis
        \item Argumentasi
        \item Penegasan ulang
        \item Koda
        \item Orientasi konflik
    \end{enumerate}

\item Kalimat di bawah ini yang memuat pernyataan berupa fakta autentik untuk teks eksposisi adalah \dots
    \begin{enumerate}[label=\alph*.]
        \item Menurut saya, kebersihan kelas ini sudah cukup baik dan perlu dipertahankan.
        \item Sebaiknya seluruh siswa membawa bekal makanan sendiri dari rumah masing-masing.
        \item Berdasarkan data resmi, luas kawasan hutan tropis Indonesia menduduki peringkat ketiga di dunia.
        \item Mungkin esok hari cuaca akan cerah sehingga kita bisa berolahraga di lapangan bebas.
        \item Menanam pohon dinilai sangat menyenangkan sekali jika dilakukan bersama keluarga.
    \end{enumerate}

\item Proses penambahan imbuhan pada suatu kata dasar yang dapat mengubah kelas kata disebut dengan proses \dots
    \begin{enumerate}[label=\alph*.]
        \item Reduplikasi
        \item Afiksasi
        \item Komposisi
        \item Akronimisasi
        \item Derivasi fiktif
    \end{enumerate}

\item Di antara kata-kata berikut, yang menggunakan konfiks (imbuhan gabung awalan-akhiran secara serentak) adalah \dots
    \begin{enumerate}[label=\alph*.]
        \item Membaca
        \item Keindahan
        \item Tulisan
        \item Makanan
        \item Telunjuk
    \end{enumerate}

\item Kata yang terbentuk melalui proses infiksasi (penyisipan imbuhan di tengah kata dasar) di bawah ini adalah \dots
    \begin{enumerate}[label=\alph*.]
        \item Bergetar
        \item Gemetar
        \item Gemetaran
        \item Menggetarkan
        \item Getaran
    \end{enumerate}

\item Pasangan kata berikut ini yang semuanya mengalami proses infiksasi dengan benar adalah \dots
    \begin{enumerate}[label=\alph*.]
        \item \textit{Gerigi} dan \textit{seruling}
        \item \textit{Keadilan} dan \textit{persatuan}
        \item \textit{Melompat} dan \textit{terjatuh}
        \item \textit{Makanan} dan \textit{minuman}
        \item \textit{Pendaftaran} dan \textit{perumahan}
    \end{enumerate}

\item Cermati kutipan teks eksposisi berikut:\\
\textit{"Pendidikan karakter sangat penting diterapkan sejak dini untuk membangun generasi bangsa yang bermoral dan berintegritas tinggi."}\\
Kutipan di atas di dalam struktur teks eksposisi menempati bagian \dots
    \begin{enumerate}[label=\alph*.]
        \item Tesis
        \item Argumentasi data
        \item Contoh kasus
        \item Penegasan ulang
        \item Daftar pustaka
    \end{enumerate}

\item Di antara kalimat-kalimat di bawah ini, yang merupakan pernyataan opini adalah \dots
    \begin{enumerate}[label=\alph*.]
        \item Candi Prambanan terletak di wilayah perbatasan Yogyakarta dan Jawa Tengah.
        \item Pemerintah seharusnya mengalokasikan anggaran lebih besar untuk sektor pertanian.
        \item Oksigen merupakan gas yang sangat dibutuhkan manusia untuk bernapas.
        \item Lagu Indonesia Raya diciptakan oleh Wage Rudolf Supratman.
        \item Air membeku pada suhu nol derajat Celsius dalam tekanan normal.
    \end{enumerate}

\item Bagian penegasan ulang dalam teks eksposisi biasanya memuat hal berikut, yaitu \dots
    \begin{enumerate}[label=\alph*.]
        \item Pertentangan baru yang memicu konflik antarpembaca.
        \item Penjelasan detail mengenai metode penelitian laboratorium ilmiah.
        \item Kesimpulan akhir yang menegaskan kembali tesis awal disertai saran/rekomendasi.
        \item Deretan angka-angka statistik tanpa kesimpulan naratif.
        \item Pengenalan tokoh utama dan latar waktu terjadinya peristiwa.
    \end{enumerate}

\item Kerangka argumentasi dalam teks eksposisi yang baik harus didasarkan pada \dots
    \begin{enumerate}[label=\alph*.]
        \item Tebakan dan asumsi subjektif tanpa bukti konkret.
        \item Data, fakta akurat, hasil penelitian, atau pendapat ahli yang valid.
        \item Cerita misteri rakyat yang berkembang di masyarakat kuno.
        \item Emosi dan amarah penulis saat menghadapi suatu masalah fiksi.
        \item Rima akhir puisi yang nadanya terdengar indah saat dibaca.
    \end{enumerate}

\item Manakah yang termasuk kalimat fakta di bawah ini? \dots
    \begin{enumerate}[label=\alph*.]
        \item Belajar bahasa Indonesia itu rasanya sangat sulit sekali bagi semua siswa.
        \item SMK Plus Pratama Adi merupakan institusi pendidikan tingkat menengah kejuruan.
        \item Masa depan seorang anak ditentukan dari seberapa mahal merk gawai yang dimilikinya.
        \item Kemungkinan besar hujan badai akan melanda kota kita sepanjang minggu depan.
        \item Menulis teks eksposisi jauh lebih mudah daripada menggambar grafik fungsi kuadrat.
    \end{enumerate}

\item Kata \textit{"keberhasilan"} diturunkan dari kata dasar \dots melalui proses konfiksasi \dots
    \begin{enumerate}[label=\alph*.]
        \item \textit{Berhasil}, konfiks ke-an
        \item \textit{Hasil}, konfiks ke-an
        \item \textit{Berhasil}, sufiks -an
        \item \textit{Hasil}, prefiks ke-
        \item \textit{Sil}, konfiks keber-an
    \end{enumerate}

\item Cermati kutipan akhir eksposisi berikut:\\
\textit{"Dengan demikian, menjaga kebersihan lingkungan sekolah merupakan tanggung jawab bersama demi kenyamanan proses belajar-mengajar. Oleh karena itu, sinergi antara siswa dan guru sangat diperlukan."}\\
Kutipan di atas merupakan bagian dari struktur \dots
    \begin{enumerate}[label=\alph*.]
        \item Tesis awal
        \item Orientasi latar
        \item Deskripsi anatomi objek
        \item Penegasan ulang (kesimpulan dan saran)
        \item Argumentasi penguat data
    \end{enumerate}

\end{enumerate}
\end{document}